\chapter{Evaluation Metrics and Result Interpretation}
\label{chap:ch16}

\section{The Multidimensional Challenge of Visual Evaluation}
How do we mathematically determine whether one reconstructed satellite image is "better" than another?

In remote sensing, human visual inspection is subjective and cannot scale to millions of square kilometers. We rely on quantitative mathematical metrics. However, blindly relying on a single metric (such as PSNR) can lead to catastrophic mistakes.

\section{Core Mathematical Metrics}

\subsection{Peak Signal-to-Noise Ratio (PSNR)}
PSNR measures the ratio between the maximum possible power of a valid physical signal and the corrupting power of reconstruction noise. Given target $\mathbf{Y}^*$ and prediction $\hat{\mathbf{Y}}$ bounded in $[0, 1]$ ($\text{MAX}_I = 1.0$):
\begin{equation}
\text{MSE} = \frac{1}{C \cdot H \cdot W} \sum_{c=1}^C \sum_{h=1}^H \sum_{w=1}^W \big( \hat{\mathbf{Y}}_{c,h,w} - \mathbf{Y}^*_{c,h,w} \big)^2
\end{equation}
\begin{equation}
\text{PSNR} = 10 \log_{10}\left( \frac{1.0^2}{\text{MSE}} \right) = -10 \log_{10}(\text{MSE}) \quad (\text{measured in decibels, dB})
\end{equation}
\begin{itemize}
    \item A higher PSNR indicates lower numerical pixel variance.
    \item A reconstructed patch with $\text{MSE} = 0.001$ yields $\text{PSNR} = -10 \log_{10}(0.001) = \mathbf{30.0\text{ dB}}$.
    \item Typical acceptable cloud removal quality is $\ge 28\text{--}32\text{ dB}$.
\end{itemize}

\subsection{Structural Similarity Index Measure (SSIM)}
Introduced by Wang et al.~\cite{wang2004ssim}, SSIM evaluates image degradation not as independent pixel errors, but as perceived changes in structural information, luminance, and contrast:
\begin{equation}
\text{SSIM}(\mathbf{x}, \mathbf{y}) = \frac{(2\mu_x \mu_y + C_1)(2\sigma_{xy} + C_2)}{(\mu_x^2 + \mu_y^2 + C_1)(\sigma_x^2 + \sigma_y^2 + C_2)}
\end{equation}
where $\mu$ denotes local mean, $\sigma$ denotes local variance, and $\sigma_{xy}$ is cross-covariance.
\begin{itemize}
    \item Range: $[-1, +1]$, where $+1.0$ indicates identical structural textures.
    \item Excellent reconstructions achieve $\text{SSIM} \ge 0.85\text{--}0.90$.
\end{itemize}

\subsection{Detail Ratio: High-Frequency Energy Retention}
Standard PSNR and SSIM are notoriously blind to blurriness. To quantify over-smoothing explicitly, our project introduced the \textbf{Detail Ratio}:
\begin{equation}
\text{Detail Ratio} = \frac{\mathbb{E}\big[ |\nabla^2 \hat{\mathbf{Y}}|^2 \big]}{\mathbb{E}\big[ |\nabla^2 \mathbf{Y}^*|^2 \big]}
\end{equation}
where $\nabla^2$ is the $3 \times 3$ discrete Laplacian operator measuring edge energy:
\[
\nabla^2 = \begin{bmatrix} 0 & 1 & 0 \\ 1 & -4 & 1 \\ 0 & 1 & 0 \end{bmatrix}
\]
\begin{itemize}
    \item A value of $1.00$ ($100\%$) indicates that the reconstruction preserves the exact high-frequency texture energy of the ground-truth clear imagery.
    \item Stage 1 baseline models scored only $\mathbf{0.23\text{--}0.47}$ ($53\%\text{--}77\%$ texture energy lost to blurring!).
    \item Our Stage 4 GAN model achieves $\mathbf{0.8576}$ ($85.8\%$ texture fidelity!).
\end{itemize}

\section{The Illusion of Global Metrics: Why Cloud Stratification is Mandatory}
Suppose a test patch has a tiny cloud covering only $5\%$ of its area ($\mathbf{M} = 1$ on $5\%$ of pixels, $0$ on $95\%$).
Because our model uses soft-mask blending:
\[
\hat{\mathbf{Y}} = (1 - \mathbf{M})\mathbf{X}_{S2} + \mathbf{M}\mathbf{P}_{\text{dir}}
\]
On $95\%$ of the pixels, $\hat{\mathbf{Y}} = \mathbf{X}_{S2} = \mathbf{Y}^*$, so the error on $95\%$ of the image is literally \textbf{zero}!
Even if the neural network completely hallucinates or fills the cloudy $5\%$ with flat grey, the global PSNR will be over $38\text{ dB}$!

\begin{warningbox}{The Necessity of Stratified Evaluation}
Reporting a single global average PSNR across an entire test set hides whether a model can actually reconstruct thick clouds. In our evaluation protocol, we enforce \textbf{cloud-coverage stratification}:
\begin{enumerate}
    \item \textbf{Heavy Cloud Strata ($70\%\text{--}100\%$ cloud coverage):} The ultimate stress test; tests true radar-to-optical reconstruction capability.
    \item \textbf{Moderate Cloud Strata ($40\%\text{--}70\%$ cloud coverage):} Tests reconstruction and contextual border integration.
    \item \textbf{Light Cloud / Haze Strata ($0\%\text{--}40\%$ cloud coverage):} Tests thin haze removal and background preservation.
\end{enumerate}
\end{warningbox}

\section{The Perception-Distortion Trade-Off}
Why does our Stage 4 GAN model have a slightly lower PSNR ($29.97\text{ dB}$) than an over-smoothed model ($30.20\text{ dB}$), yet looks dramatically sharper and more realistic to human geologists?

This phenomenon is known in computer vision as the \textbf{Perception-Distortion Trade-Off} (Blau \& Michaeli):
\begin{itemize}
    \item If a model predicts a sharp agricultural tree row shifted by just 1 pixel relative to the target, the pixel-wise difference $(\hat{y}_i - y_i)^2$ is penalized twice (once at the old location, once at the new one), driving PSNR down!
    \item If a model blurs the trees into a smooth green wash, its error at every pixel is moderate, yielding a higher PSNR!
    \item A higher PSNR can reward blurriness. True research evaluation requires evaluating PSNR alongside SSIM, Edge Correlation, and the Detail Ratio.
\end{itemize}

\section{Summary of Chapter 16}
\begin{enumerate}
    \item PSNR measures logarithmic signal-to-noise ratio; SSIM measures structural and contrast similarity.
    \item The Detail Ratio computes Laplacian high-frequency energy retention, exposing over-smoothing.
    \item Global metrics are skewed by clear-sky pixels; evaluation must be stratified by cloud coverage ($>70\%$ heavy clouds).
    \item The perception-distortion trade-off explains why perceptually sharper GAN models can have slightly lower PSNR than artificially blurred models.
\end{enumerate}

\section{Practice Exercises}
\begin{enumerate}
    \item Calculate the PSNR in decibels (dB) for an image reconstruction with Mean Squared Error $\text{MSE} = 0.0001$.
    \item A model achieves a Detail Ratio of $0.35$. What does this reveal about its physical reconstructions?
    \item Why does a 1-pixel spatial shift of a sharp edge severely reduce PSNR while barely affecting human perceptual quality?
\end{enumerate}

\subsection*{Solutions}
\begin{enumerate}
    \item $\text{PSNR} = -10 \log_{10}(0.0001) = -10 \log_{10}(10^{-4}) = -10 \times (-4) = \mathbf{40.0\text{ dB}}$.
    \item A Detail Ratio of $0.35$ means that the reconstructed image contains only $35\%$ of the high-frequency edge energy of the true clear scene; $65\%$ of the fine texture has been smoothed away into blurry gradients.
    \item Pixel-wise MSE compares values strictly at the identical coordinate $(x, y)$. If a high-contrast edge ($0 \to 1$) shifts by 1 pixel, the MSE at that pixel jumps from $0 \to (1 - 0)^2 = 1.0$. Even though the edge is structurally present, the point-wise error is maximal, heavily reducing the logarithmic PSNR score.
\end{enumerate}
